\section{Zaključak}
\label{sec:zakljucak}

Ovaj seminarski rad sustavno je prikazao ulogu dviju ključnih normi — ISO~26262 i ISO~21448 (SOTIF) — u osiguravanju sigurnosti sustava automatizirane vožnje, smjestivši ih u tehnički kontekst arhitekture i percepcije ADS-a. Analiza je pokazala da te dvije norme pokrivaju komplementarne aspekte sigurnosti: ISO~26262 upravlja rizicima koji proizlaze iz kvarova E/E sustava, dok ISO~21448 adresira opasnosti uzrokovane funkcionalnim insuficijencijama i ograničenjima performansi sustava koji rade prema specifikaciji.

Prikaz perceptivnog lanca pokazao je zašto je ova podjela nužna. Suvremeni ADS oslanja se na fuziju komplementarnih senzora — kamera, LiDAR-a i radara — od kojih nijedan ne pruža pouzdanu percepciju u svim uvjetima. Upravo okolišni uvjeti poput magle, jake kiše ili zasljepljujućeg sunca, koji degradiraju percepciju iako senzori rade ispravno, predstavljaju tipične okidajuće uvjete u smislu SOTIF-a. Time se uspostavlja izravna veza između inženjerske izvedbe percepcije i normativnog okvira sigurnosti.

Norma ISO~26262, sa svojim V-modelom razvoja, HARA metodologijom, ASIL klasifikacijom te hardverskim i procesnim mjerama, pruža robusnu i dobro uspostavljenu strukturu za upravljanje funkcijskom sigurnošću. Međutim, rastuća složenost sustava automatizirane vožnje — posebno onih temeljenih na strojnom učenju i naprednoj senzorskoj fuziji — nadilazi opseg tradicionalne funkcijske sigurnosti. Upravo tu ISO~21448 popunjava kritičnu prazninu, uvodeći koncepte operativne domene, okidajućih uvjeta, klasifikacije scenarija i iterativne validacije koji su prilagođeni inherentnoj nesigurnosti senzorsko-perceptivnih sustava.

Integracija dviju normi u koherentan razvojni proces ostaje značajan izazov, ponajprije zbog nepodudarnosti njihovih razvojnih modela (linearan V-model naspram iterativnog SOTIF ciklusa), složenosti upravljanja dvostrukim metodologijama te nedostatka konsenzusa oko kvantitativnih kriterija prihvatljivosti rezidualnog rizika. Ipak, pregledana literatura ukazuje na ostvarivost holističkog pristupa koji objedinjuje funkcijsku sigurnost i SOTIF unutar zajedničkog životnog ciklusa, dok prijedlozi poput mjerljivih dimenzija upravljivosti (prenosivost i predvidljivost) pokazuju kako se postojeća struktura ISO~26262 može proširiti za bezvozačke sustave razina SAE~4 i SAE~5, a da se pritom očuva njezina rigoroznost.

Gledajući u budućnost, normativni okvir za automatizirane vozne sustave nastavit će se razvijati prema holističkom pristupu koji uz ISO~26262 i ISO~21448 obuhvaća i scenarijsku evaluaciju (ISO~34502), verifikaciju i validaciju specifičnu za ADS (ISO/TR~4804), formalizaciju pretpostavki o ponašanju sudionika u prometu (IEEE~2846) te sveobuhvatni sigurnosni slučaj (UL~4600). Za inženjere i istraživače, razumijevanje komplementarnosti ISO~26262 i ISO~21448 — kao i njihove veze s arhitekturom percepcije — predstavlja neophodan temelj za suočavanje s izazovima koje donose vozila visokih razina automatizacije.
